% ===================================================================
%  اللوحة رقم ١٤ — الشارع. — التجّار.
%  Traduction 2026 d'apres texto/io, controlee sur texto/fr, texto/en
%  et texto/es. Consigne : voir texto/ar/10-lawha-01.tex.
% ===================================================================

%%K t14-apar-1 apar
\VUsaut{4.0mm}
\VUtitre{15.0pt}{60}{اللوحة رقم ١٤}
\VUsaut{3.0mm}

%%K t14-tit-1 sub
\VUpk{11.4pt}{40}{الشارع. --- التجّار.}
\VUsaut{3.0mm}

%%K t14-01-1 p
1. --- جاء صديقي يوحنا، وهو يسكن الريف، ليقضي ثلاثة أيام عند أهلي.
و\VUgras{حتى} الآن لم تسنح له قط فرصة رؤية مدينة كبيرة، فنحن نقضي أيامنا
كلها في التجوال وأنا أشرح له ما لا يعرف.

%%K t14-02-1 p
2. --- ذهبنا أولًا لنرى \VUgras{المرصد} \textsuperscript{(1)} فأعجبه أشد الإعجاب.
وفي عودتنا من \VUgras{الشارع} \textsuperscript{(3)} الذي فيه \VUgras{الحمّام
التركي} \textsuperscript{(2)}، صادفنا \VUgras{جنازة} \textsuperscript{(4)}. وفي مقدمة
\VUgras{موكب الجنازة} يمشي \VUgras{رجال الدين} أمام \VUgras{عربة الموتى}
التي وُضع فيها \VUgras{النعش}. ويتبع أصدقاء كثيرون الأسرة في
\VUgras{حِدادها}.

%%K t14-02-2 p
ولما مرّ الموكب عبرنا الشارع أمام \VUgras{الحانة} \textsuperscript{(5)} الكبيرة،
حيث كان \VUgras{زبائن} كثيرون يشربون \VUgras{الجعة} على الشرفة في ظل
\VUgras{المظلة} \textsuperscript{(6)} الزجاجية.

%%K t14-03-1 p
3. --- وحيّر يوحنا \VUgras{شعار النبالة} الموضوع بين دكان \VUgras{تاجر
الخمور والمشروبات} \textsuperscript{(7)} ودكان \VUgras{بائع النوتات}
\textsuperscript{(10)}. فسألني ما معناه؛ فقلت له إنه يدل على \VUgras{القنصلية}
\textsuperscript{(8)} الإنجليزية، وأشرت له إلى \VUgras{القنصل} \textsuperscript{(9)}
الواقف في الشرفة.

%%K t14-tit-2 sub
\VUpk{10.2pt}{40}{النظر في الواجهات.}
\VUsaut{3.0mm}

%%K t14-04-1 p
4. --- ووقفنا بطبيعة الحال أمام معرض \VUgras{الآلات الموسيقية}
و\VUgras{الأرغنات الصغيرة} و\VUgras{الأرغنات} و\VUgras{البيانوهات}؛ ونظرنا
في الأغلفة المصوّرة لـ\VUgras{النوتات} و\VUgras{قطع} البيانو، وأعجبتنا
\VUgras{القيثارة} الخشبية المذهّبة المتخذة لافتةً.

%%K t14-05-1 p
5. --- وسُحب الغبار التي أثارها \VUgras{كنّاسو الشوارع} \textsuperscript{(11)}
و\VUgras{آلة الكنس} \textsuperscript{(12)} اضطرتنا إلى الإسراع أمام
\VUgras{أطقم الأثاث} الجميلة عند \VUgras{تاجر الأثاث} \textsuperscript{(13)}
وأمام معرض \VUgras{المدافئ} و\VUgras{المصابيح} عند \VUgras{بائع الخردوات}
\textsuperscript{(14)}.

%%K t14-06-1 p
6. --- وفي ذلك الموضع اضطررنا فعلًا إلى ترك الرصيف، فقد سدّته الطرود
التي كانت تُنزل من \VUgras{عربة نقل الأثاث} \textsuperscript{(16)} لتُحمل إلى
\VUgras{دكان} \textsuperscript{(15)} أُجّر لتوّه.

%%K t14-06-2 p
وحال ذلك دون رؤيتنا العربات الجميلة عند \VUgras{صانع العربات}
\textsuperscript{(17)}، لكنه لم يمنعنا من الوقوف لنرى \VUgras{المعبّئ}
\textsuperscript{(19)} يصنع صندوقًا لجاره \VUgras{تاجر الدراجات والسيارات}
\textsuperscript{(18)}. ثم عدنا إلى البيت، إذ كان الوقت قد تأخر بعض الشيء.

%%K t14-tit-3 sub
\VUpk{10.2pt}{40}{جولة في المدينة.}
\VUsaut{3.0mm}

%%K t14-07-1 p
7. --- وفي الغد اقترحت أمي أن يُصوَّر يوحنا، وطلبت مني أن أرافقه إلى
\VUgras{المصوّر} \textsuperscript{(20)}. وبعد الغداء ذهبنا إلى \VUgras{مرسم}
الفنان حيث انتظرنا طويلًا. ولنقطع الوقت نظرنا في \VUgras{الصور}
\textsuperscript{(22)} المعلّقة على الجدار.

%%K t14-07-2 p
ولما جاء دورنا لم يطل الأمر، وما إن فرغ صديقي من الوقوف أمام
\VUgras{آلة التصوير} \textsuperscript{(21)} حتى نزلنا.

%%K t14-08-1 p
8. --- وفي الطابق الأول رأينا، من باب \VUgras{الحلّاق} \textsuperscript{(23)}
المفتوح، مساعده يقصّ شعر زبون.

%%K t14-08-2 p
وعدنا إلى الشارع فمررنا بـ\VUgras{دكان التبغ} \textsuperscript{(24)}. وقد أضحك
يوحنا \VUgras{الفانوس} \textsuperscript{(25)} الأحمر و\VUgras{الغليون الكبير}
المتخذ \VUgras{لافتةً} \textsuperscript{(26)} حتى اصطدم بشيخين يتحادثان على
\VUgras{نشقة سعوط} \textsuperscript{(27)}.

%%K t14-tit-4 sub
\VUpk{10.2pt}{40}{دكان الحلوى.}
\VUsaut{3.0mm}

%%K t14-09-1 p
9. --- و\VUgras{الأوراق النقدية} و\VUgras{القطع المعدنية} من كل بلد
المعروضة في \VUgras{واجهة} \VUgras{الصرّاف} \textsuperscript{(29)} شدّت انتباهنا
لحظة، لكن الوقت كان وقت العصرية فدخلنا \VUgras{دكان الحلواني}
\textsuperscript{(31)} دون مزيد وقوف.

%%K t14-10-1 p
10. --- ولما رأينا كل \VUgras{الطيبات} التي في الدكان --- \VUgras{كعك
وخبز زنجبيل وبسكويت} و\VUgras{حلوى} من كل نوع --- عزّ علينا الاختيار. وفي
الآخر استقرّ رأي يوحنا على \VUgras{كعك القشدة}، وأما أنا فلم آخذ إلا
\VUgras{فطيرة كرز} صغيرة، فالحلويات \VUgras{تسبّب لي وجع الأسنان} ولا رغبة
لي البتة في زيارة \VUgras{طبيب الأسنان} \textsuperscript{(30)} كثيرًا.

%%K t14-tit-5 sub
\VUpk{10.2pt}{40}{الكتابة على الآلة.}
\VUsaut{3.0mm}

%%K t14-11-1 p
11. --- ولأُري صديقي \VUgras{آلة كاتبة} كان قد سمع بها ولم يرها قط، دخلنا
\VUgras{مكتب} \textsuperscript{(32)} \VUgras{ابن عمي} \textsuperscript{(34)}. فوجدناه
يملي رسائله على \VUgras{سكرتيرته} \textsuperscript{(35)}، وهي \VUgras{مختزلة}. وما
تكاد رسالة تنتهي حتى ينسخها \VUgras{كاتب على الآلة} \textsuperscript{(33)}. ولأننا
لم نرد قطع عمل قريبي، لم نمكث عنده إلا دقائق.

%%K t14-12-1 p
12. --- وتحت مكتب ابن عمي يسكن \VUgras{ساعاتي وجواهرجي} \textsuperscript{(37)}
كبير، وهو صائغ فضة أيضًا. وفي دكانه البديع \VUgras{ساعات حائط وساعات رفّ
وساعات جيب} كثيرة، و\VUgras{سلاسل}، و\VUgras{جواهر} ثمينة، مثل
\VUgras{عقود اللؤلؤ} و\VUgras{الأساور} الذهبية و\VUgras{الأقراط}
و\VUgras{الخواتم} المرصّعة بـ\VUgras{الأحجار الكريمة} و\VUgras{الألماس}.
وإحدى الواجهات مخصّصة كلها لـ\VUgras{الفضّيات} وفيها مجموعة كبيرة من
\VUgras{أدوات المائدة} الفضية.

%%K t14-13-1 p
13. --- وعند انعطافي عند زاوية الشارع لاحظت أن \VUgras{اللوحة}
\textsuperscript{(36)} التي إلى جانب \VUgras{رقم} البيت قد بُدّلت. وكنت سأقول ذلك
بصوت عالٍ حين سُمع الصوت المرح لفرقة عسكرية. فأخذنا نجري نحو الفوج، دون
أن نفكر أنه سيمرّ أمام دكاني \VUgras{القبّعاتي} \textsuperscript{(39)}
و\VUgras{القفّازاتي} \textsuperscript{(40)}، وأننا كنا سنكون في موضع مثله جودةً
لنرى استعراضه لو بقينا أمام واجهة \VUgras{النظّاراتي} \textsuperscript{(38)}
الملأى بـ\VUgras{النظارات الأنفية والنظارات} و\VUgras{مناظير المسرح}.

%%K t14-tit-6 sub
\VUpk{10.2pt}{40}{الاستعراض.}
\VUsaut{3.0mm}

%%K t14-14-1 p
14. --- في مقدمة \VUgras{الفوج} \textsuperscript{(41)} كان يسير \VUgras{قائد
الطبّالين} \textsuperscript{(42)}، يتبعه \VUgras{الطبّالون} \textsuperscript{(43)}
و\VUgras{البوّاقون} \textsuperscript{(44)} و\VUgras{الفرقة} كلها. و\VUgras{الجنرال}
\textsuperscript{(45)}، وكان في الساحة مع مرافقه، قد وقف قرب \VUgras{نافورة}
\textsuperscript{(49)} \VUgras{الحوض} \textsuperscript{(48)} ليرى \VUgras{الاستعراض}.

%%K t14-14-2 p
و\VUgras{ضباط الأركان} \textsuperscript{(46)} و\VUgras{العقيد} و\VUgras{المقدّم}
يحيّونه بسيوفهم. و\VUgras{الرواد} في مقدمة \VUgras{كتائبهم} \textsuperscript{(47)}،
وكل \VUgras{نقيب} يقود \VUgras{سريّته}، بينما يقف \VUgras{الملازمون}
و\VUgras{الملازمون الثواني} و\VUgras{ضباط الصف} في مواضعهم المعتادة إلى
جانب رجالهم.

%%K t14-15-1 p
15. --- وقد تجمّع كثير من المارّة عند \VUgras{مفترق الطرق} \textsuperscript{(50)}؛
وخرج أصحاب الدكاكين --- \VUgras{بائع المظلات} \textsuperscript{(51)}،
و\VUgras{الصبّاغ} \textsuperscript{(53)}، و\VUgras{بائع الأسلحة} \textsuperscript{(54)} ---
ليروا \VUgras{الجنود}. واصطفّت المركبات كلها --- \VUgras{عربات الأجرة}
\textsuperscript{(52)}، و\VUgras{العربات الخاصة} \textsuperscript{(57)}، بل و\VUgras{صهريج
الرشّ} \textsuperscript{(56)} --- على طول الرصيف.

%%K t14-tit-7 sub
\VUpk{10.2pt}{40}{مشاهد من الشارع.}
\VUsaut{3.0mm}

%%K t14-15-2 p
و\VUgras{مندوب تجاري} \textsuperscript{(55)} يحمل \VUgras{حقائب عيّناته} في يده
عبر الشارع مسرعًا، والتفت الجالسون على \VUgras{الطابق العلوي} \textsuperscript{(59)}
من \VUgras{العربة العامة} \textsuperscript{(58)} ليتمتعوا بالمنظر. واضطر
\VUgras{مغنٍّ جوّال} \textsuperscript{(60)} مسكين و\VUgras{أرغنه} \textsuperscript{(61)}
معه إلى قطع \VUgras{أغنيته}، فقد غطّى دويّ الطبول على صوته.

%%K t14-16-1 p
16. --- ولما بلغ الفوج أمام \VUgras{دكان الأقمشة} \textsuperscript{(64)}، تركت
\VUgras{الخياطات} \textsuperscript{(63)} و\VUgras{الخياطون} \textsuperscript{(62)}
\VUgras{آلات الخياطة} وعملهم ليطلّوا من النافذة. واضطر \VUgras{صاحب الدكان}
\textsuperscript{(65)}، وقد وضع لتوّه \VUgras{بذلات} \textsuperscript{(66)} جديدة في
\VUgras{واجهته}، إلى نقل لفّات \VUgras{الجوخ} \textsuperscript{(67)}
و\VUgras{الكتّان} المبسوطة على الرصيف قرب \VUgras{فوهة البالوعة}
\textsuperscript{(68)}، مخافة أن يقلبها الحشد؛ حتى \VUgras{بائع متجول}
\textsuperscript{(69)} عجوز، كانت بوّابة تساومه على جوارب، اضطر إلى الدخول في
مدخل بيت ليتم بيعه.

%%K t14-16-2 p
ورافقنا الفوج إلى الثكنة \VUgras{ثم} عدنا إلى البيت وقت العشاء، وقد سررنا
بيومنا سرورًا شديدًا.

%%K t14-tit-8 sub
\VUpk{10.2pt}{40}{حرف ومهن شتى.}
\VUsaut{3.0mm}

%%K t14-17-1 p
17. --- وفي الغد عرض علينا أبي أن يأخذنا في جولة في مطبعة الجريدة
المحلية. فقبلنا بحماس.

%%K t14-17-2 p
و\VUgras{الطابع} هو أيضًا \VUgras{بائع كتب} \textsuperscript{(70)}
و\VUgras{ناشر} و\VUgras{بائع قرطاسية} و\VUgras{مجلّد}. وقد وضع في الطابق
الأول \VUgras{هيئة تحرير} \textsuperscript{(71)} الجريدة، وكذلك \VUgras{قاعة
الحفر} حيث يعمل \VUgras{الطابعون على الحجر} \textsuperscript{(72)}
و\VUgras{الحفّارون على النحاس} \textsuperscript{(73)}، وأخيرًا قاعة الجمع حيث
\VUgras{الجامعون} \textsuperscript{(74)}. أما \VUgras{قاعة المكابس} \textsuperscript{(75)}
نفسها و\VUgras{مطابع الطبع} والآلات المختلفة فهي في الطابق الأرضي.

%%K t14-tit-9 sub
\VUpk{10.2pt}{40}{يوحنا يسأل أسئلة كثيرة جدًا.}
\VUsaut{3.0mm}

%%K t14-18-1 p
18. --- ولما خرجنا من المطبعة وقف يوحنا في حيرة شديدة أمام علبة زجاجية
صغيرة مركّبة على عمود قبالة \VUgras{دكان الخردوات} \textsuperscript{(77)}، وسأل
ما فائدتها. فشرح له أبي أنها \VUgras{جهاز إنذار حريق} \textsuperscript{(76)}. بل
إن كل شيء في المدينة كان عجبًا لصديقي: من \VUgras{نافورة الشرب}
\textsuperscript{(78)} البسيطة و\VUgras{الدراجة الثلاثية للتوصيل} \textsuperscript{(80)}
الخفيفة يركبها \VUgras{صبي} \textsuperscript{(79)} في زيّ رسمي، إلى \VUgras{عربة
البريد} \textsuperscript{(81)}.

%%K t14-19-1 p
19. --- وبينما تركنا أبي لحظة ليكلّم \VUgras{جابي الضرائب} \textsuperscript{(83)}
الذي وقف يتحدث مع \VUgras{محامٍ} \textsuperscript{(82)} و\VUgras{موثّق}
\textsuperscript{(84)}، تسلّينا بمتابعة حركة الشارع. فمرّ أمامنا أشد الناس
تنوعًا: \VUgras{صبي حلواني} \textsuperscript{(85)} يحمل في سلة \VUgras{فطيرة}
بديعة، جاء خلف \VUgras{بائع ثياب مستعملة} \textsuperscript{(86)}؛ و\VUgras{مشعل
الفوانيس} \textsuperscript{(87)} يتحدث مع \VUgras{عامل غاز} \textsuperscript{(88)}؛
و\VUgras{راهبة} \textsuperscript{(89)} تمسك بيد يتيمة صغيرة عائدة إلى ديرها.

%%K t14-tit-10 sub
\VUpk{10.2pt}{40}{في طريق العودة.}
\VUsaut{3.0mm}

%%K t14-20-1 p
20. --- وفي اللحظة التي عاد فيها أبي إلينا، انفجر يوحنا ضاحكًا من الوجوه
المسخّمة والهيئة الغريبة لـ\VUgras{كاسحَي مداخن} \textsuperscript{(91)} أسودين
من السخام.

%%K t14-21-1 p
21. --- وأردت أن أشتري لأمي نبتة من \VUgras{بائعة الزهور} \textsuperscript{(92)}،
لكن أبي لاحظ أنها ستكون ثقيلة الحمل وأن الأولى شراء \VUgras{برتقال}.
فتقدمنا إلى \VUgras{صاحبة البسطة} \textsuperscript{(94)}، ولما فرغت
\VUgras{المربّية} \textsuperscript{(93)}، وكانت تختار بعضه لصغيرها، من شرائها،
جاء دورنا.

%%K t14-22-1 p
22. --- وفي طريق العودة التقينا بعدة معارف: صديقة قديمة لأسرتي متكئة على
ذراع \VUgras{وصيفتها} \textsuperscript{(95)}، و\VUgras{مهندس} \textsuperscript{(96)}
الطرق والجسور، وإحدى بنات عمي مع \VUgras{مربّيتها} \textsuperscript{(98)}.

%%K t14-22-2 p
ثم مررنا بـ\VUgras{مصمّمة القبعات} \textsuperscript{(97)} التي تخيط لأمي، وبـ\VUgras{راهبين}
\textsuperscript{(99)} غريبين عن المدينة، أقبلا يسألان أبي عن
الطريق إلى المحطة.
\VUsaut{6.0mm}
\VUfilet{22mm}
